\documentclass[12pt,a4paper]{article}

% ------------------ PAQUETES ------------------
\usepackage[utf8]{inputenc}
\usepackage[T1]{fontenc}
\usepackage[spanish,es-tabla]{babel} 
\usepackage{geometry}
\usepackage{setspace}
\usepackage{hyperref}
\usepackage{pgfgantt}
\usepackage{float}
\usepackage{enumitem} 
\usepackage{colortbl} % Para colorear celdas
\usepackage{xcolor}   
\usepackage{array}    % Nuevo paquete para mejorar un poco las tablas

\geometry{margin=2.5cm}
\setstretch{1.2}      
\definecolor{headergray}{gray}{0.9}

% ------------------ TÍTULO ------------------
\title{
    \vspace{-2cm}
    \Large \textbf{Ejecución de la planificación y Cronograma del TFG}\\[10pt]
    \large Diseño y análisis de un agente LLM autónomo para auditorías de ciberseguridad
}

\author{\textbf{Alejandro Barrena Millán}}
\date{3 de diciembre de 2025}

\begin{document}

\maketitle
\thispagestyle{empty} 
\hrule
\vspace{1em}

\section{Introducción}

El presente documento detalla la planificación temporal y la metodología de trabajo establecida para el desarrollo del Trabajo de Fin de Grado (TFG) titulado provisionalmente \textit{Diseño y análisis de un agente LLM autónomo para auditorías de ciberseguridad en entornos aislados}.

El objetivo es definir un calendario estructurado que garantice el cumplimiento de los hitos de desarrollo y la entrega de la memoria final en los plazos previstos.

% ============================================================
\section{Metodología: Scrum Adaptado}

El desarrollo del proyecto se regirá por un enfoque ágil basado en \textit{Scrum}, adaptado a las necesidades de un trabajo individual. Este marco de trabajo permite iterar sobre el proyecto, asegurar entregables funcionales en cada etapa y adaptar la planificación ante posibles imprevistos técnicos.

El ciclo de vida se divide en \textbf{Sprints} (ciclos de trabajo) con los siguientes eventos:
\begin{itemize}
    \item \textbf{Sprint Planning:} Definición de tareas y alcance al inicio de cada fase.
    \item \textbf{Desarrollo e Implementación:} Ejecución técnica de las tareas.
    \item \textbf{Sprint Review \& Retrospective:} Validación del entregable y ajuste de la metodología para el siguiente ciclo.
\end{itemize}

% ============================================================
\section{Planificación de Sprints y Tareas}

A continuación, se describe el desglose de tareas organizado por Sprints. 

\begin{description}[style=nextline, leftmargin=1.5cm, font=\bfseries]

    \item[Sprint 1: Definición y Organización (Semana 1)]
    Establecimiento de las bases del proyecto.
    \begin{itemize}[leftmargin=0.5cm, label=$\circ$]
        \item Definición detallada de requisitos funcionales y no funcionales.
        \item Selección y configuración de herramientas de gestión (e.g., Jira/Trello, GitHub).
        \item Investigación del estado del arte preliminar.
    \end{itemize}

    \item[Sprint 2: Entorno de Pruebas (Semanas 2 - 3)]
    Configuración de la infraestructura segura.
    \begin{itemize}[leftmargin=0.5cm, label=$\circ$]
        \item Despliegue del entorno de terminal en sandbox (Docker/VM).
        \item Desarrollo de la API de interfaz para la comunicación con el entorno.
        \item Pruebas de conectividad y seguridad del entorno aislado.
    \end{itemize}

    \item[Sprint 3: Integración del LLM (Semanas 4 - 6)]
    Conexión del motor de inteligencia artificial con el entorno.
    \begin{itemize}[leftmargin=0.5cm, label=$\circ$]
        \item Selección y configuración del modelo LLM (local o API externa).
        \item Desarrollo del módulo de interpretación de comandos.
        \item Prototipado inicial: el agente puede enviar comandos simples y recibir *outputs*.
    \end{itemize}

    \item[Sprint 4: Desarrollo del MVP (Semanas 7 - 8)]
    Implementación de la lógica de auditoría.
    \begin{itemize}[leftmargin=0.5cm, label=$\circ$]
        \item Programación de agentes autónomos con capacidad de decisión.
        \item Implementación de flujos de auditoría básica (e.g., escaneo de puertos, identificación de servicios).
        \item Validación del Producto Mínimo Viable (MVP).
    \end{itemize}

    \item[Sprint 5: Validación y Pruebas (Semana 9)]
    Consolidación del sistema.
    \begin{itemize}[leftmargin=0.5cm, label=$\circ$]
        \item Ejecución de auditorías *end-to-end* automatizadas.
        \item Análisis de falsos positivos y refinamiento de *prompts*.
        \item Pruebas de estrés y estabilidad.
    \end{itemize}

    \item[Sprint 6: Documentación y Memoria (Semanas 10 - 11)]
    Redacción académica.
    \begin{itemize}[leftmargin=0.5cm, label=$\circ$]
        \item Redacción de la memoria del TFG.
        \item Generación de gráficos, tablas de resultados y bibliografía.
        \item Revisión de formato y estilo.
    \end{itemize}

    \item[Sprint 7: Defensa (Semana 12)]
    Preparación final.
    \begin{itemize}[leftmargin=0.5cm, label=$\circ$]
        \item Diseño de diapositivas de presentación.
        \item Ensayo del discurso y preparación de la demo (si aplica).
    \end{itemize}

\end{description}

% ============================================================
\section{Cronograma Visual y Calendario}

Para facilitar el seguimiento del proyecto, se establece la correspondencia entre las semanas de trabajo y las fechas de calendario real, seguido de la representación visual mediante diagrama de Gantt.

\subsection{Calendario Detallado}

\begin{table}[H]
\centering
\renewcommand{\arraystretch}{1.3} % Más espacio entre filas
\begin{tabular}{|c|l|l|}
\hline
\rowcolor{headergray} \textbf{Semana} & \textbf{Fecha de Inicio (Lunes)} & \textbf{Fase / Sprint Activo} \\ \hline
1  & 17 de Noviembre, 2025 & S1: Definición y Organización \\ \hline
2  & 24 de Noviembre, 2025 & S2: Entorno Sandbox (Inicio) \\ \hline
3  & 01 de Diciembre, 2025 & S2: Entorno Sandbox (Cierre) \\ \hline
4  & 08 de Diciembre, 2025 & S3: Integración LLM (Inicio) \\ \hline
5  & 15 de Diciembre, 2025 & S3: Integración LLM (Desarrollo) \\ \hline
6  & 22 de Diciembre, 2025 & S3: Integración LLM (Cierre) \\ \hline
7  & 29 de Diciembre, 2025 & S4: Desarrollo MVP (Inicio) \\ \hline
8  & 05 de Enero, 2026     & S4: Desarrollo MVP (Cierre) \\ \hline
9  & 12 de Enero, 2026     & S5: Validación y Pruebas \\ \hline
10 & 19 de Enero, 2026     & S6: Documentación (Borrador) \\ \hline
11 & 26 de Enero, 2026     & S6: Documentación (Revisión) \\ \hline
12 & 02 de Febrero, 2026   & S7: Defensa y Presentación \\ \hline
\end{tabular}
\caption{Correspondencia de semanas con fechas naturales.}
\end{table}

\subsection{Diagrama de Gantt}

El diagrama ilustra la distribución temporal relativa de los Sprints a lo largo de las 12 semanas.

\begin{figure}[H]
\centering
\begin{ganttchart}[
    hgrid,
    vgrid,
    x unit=1.1cm,
    y unit chart=0.6cm,
    bar height=0.6,
    group right shift=0,
    group top shift=0.7,
    group height=.3,
    title label font=\bfseries\footnotesize,
    bar label font=\footnotesize
]{1}{12}

% Títulos de las semanas
\gantttitle{Semanas de Desarrollo}{12} \\
\gantttitlelist{1,...,12}{1} \\

% Barras del cronograma (rojo = ejecutado, azul = planificado)
\ganttbar[bar/.append style={fill=red!60, draw=none}]{S1: Organización}{1}{1} \\
\ganttbar[bar/.append style={fill=red!60, draw=none}]{S2: Entorno Sandbox}{2}{3} \\
\ganttbar[bar/.append style={fill=blue!40, draw=none}]{S3: Integración LLM (plan)}{4}{6} \\
\ganttbar[bar/.append style={fill=blue!40, draw=none}]{S4: MVP Agente}{7}{8} \\
\ganttbar[bar/.append style={fill=blue!40, draw=none}]{S5: Validación}{9}{9} \\
\ganttbar[bar/.append style={fill=blue!40, draw=none}]{S6: Memoria}{10}{11} \\
\ganttbar[bar/.append style={fill=blue!40, draw=none}]{S7: Defensa}{12}{12}     

% Relaciones 
\ganttlink{elem0}{elem1}
\ganttlink{elem1}{elem2}
\ganttlink{elem2}{elem3}
\ganttlink{elem3}{elem4}
\ganttlink{elem4}{elem5}
\ganttlink{elem5}{elem6}

\end{ganttchart}
\caption{Cronograma general del proyecto.}

% Leyenda de colores
\vspace{0.5em}
\begin{center}
\begin{tabular}{ll}
\colorbox{red!60}{\rule{1cm}{0.4cm}} & Ejecutado (fechas reales) \\
\colorbox{blue!40}{\rule{1cm}{0.4cm}} & Planificado (cronograma inicial) \\
\end{tabular}
\end{center}

\end{figure}

% ============================================================
\section{Conclusión}

La planificación presentada asegura un flujo de trabajo constante y controlado. La estructura modular de los Sprints permite afrontar la complejidad de integrar Modelos de Lenguaje (LLM) en entornos de ciberseguridad paso a paso, mitigando riesgos técnicos antes de la fase final de documentación.

% --------------------------------------------------
% Secciones de estado real (actualizaciones del proyecto)
% --------------------------------------------------
\newpage
\section{Estado real de los Sprints}

Las siguientes subsecciones recogen el estado real observado durante la ejecución del proyecto, con anotaciones de fechas y comentarios sobre el progreso respecto a la planificación inicial.

% Definición de colores semánticos
\definecolor{ontime}{RGB}{66, 133, 244}    % Azul - a tiempo
\definecolor{early}{RGB}{52, 168, 83}      % Verde - anticipado
\definecolor{late}{RGB}{234, 67, 53}       % Rojo - retrasado
\definecolor{inprogress}{RGB}{251, 188, 4} % Amarillo - en progreso

\subsection{Sprint 1: Definición y Organización}

\textbf{Estado:} \textcolor{ontime}{\textbf{Completado (según calendario)}}

\textbf{Fecha de cierre:} 23 de noviembre de 2025.

\textbf{Comentarios:}
\begin{itemize}
    \item Definición de requisitos funcionales y no funcionales completada.
    \item Herramientas de gestión y repositorio configurados; rama principal inicializada.
    \item Investigación y documentación preliminar realizadas.
\end{itemize}

\subsection{Sprint 2: Entorno de Pruebas}

\textbf{Estado:} \textcolor{early}{\textbf{Completado (finalizado anticipadamente)}}

\textbf{Fecha de cierre:} 02 de diciembre de 2025 \textit{(planificado: 07/12/2025)}.

\textbf{Comentarios:}
\begin{itemize}
    \item La orquestación del laboratorio y la integración mediante \texttt{docker-compose} se completaron antes de la fecha prevista.
    \item El agente API quedó operativo y se validó la comunicación con los entornos vulnerables (DVWA, Metasploitable2, OWASP Juice Shop).
    \item Se añadieron mejoras al contenedor del agente (persistencia de directorios, utilidades adicionales) y se realizaron pruebas de validación funcionales.
    \item Resultado: Sprint cerrado y entregable aceptado como base para la integración del LLM.
\end{itemize}

\subsection{Sprint 3: Integración del LLM - "El Cerebro" y Abstracción de Modelos}

\textbf{Estado:} \textcolor{ontime}{\textbf{Completado (según calendario ajustado)}}

\textbf{Fecha de inicio:} 03 de diciembre de 2025.
\textbf{Fecha de cierre:} 25 de diciembre de 2025.

\textbf{Comentarios:}
\begin{itemize}
    \item Desarrollo completo del Backend utilizando FastAPI.
    \item Implementación del sistema de Ranking ELO para evaluación de modelos.
    \item Sistema de descubrimiento automático de modelos LLM.
    \item Integración exitosa con múltiples proveedores de modelos (OpenAI, Anthropic, locales).
    \item El sprint se completó según la planificación ajustada, aprovechando el margen ganado en el Sprint 2.
\end{itemize}

\subsection{Sprint 4: Desarrollo MVP - Frontend y Mission Board}

\textbf{Estado:} \textcolor{ontime}{\textbf{Completado (según calendario)}}

\textbf{Fecha de inicio:} 26 de diciembre de 2025.
\textbf{Fecha de cierre:} 08 de enero de 2026.

\textbf{Comentarios:}
\begin{itemize}
    \item Desarrollo completo del Frontend utilizando Next.js y React.
    \item Implementación del Mission Board para gestión de tareas del agente.
    \item Desarrollo de la lógica de "anti-resignación" del agente.
    \item Integración exitosa entre Frontend y Backend.
\end{itemize}

\subsection{Reestructuración del Proyecto y Ajustes de Cronograma}
Tras la finalización del Sprint 4, se identificaron oportunidades críticas de mejora técnica que justificaron una reestructuración del cronograma original dado a que llevabamos un adelanto respecto al calendario desde la finalización del Sprint 2.
\begin{itemize}
    \item Acceso y pruebas con infraestructura Computaex de alto rendimiento, optimizando la ejecución de modelos LLM más complejos.
    \item Descarga de modelos pesados compatibles con el computaex que enriquecerán las comparativas.
    \item Implementación de un sistema de gestión de ramas de conversaciones para mejorar la persistencia y seguimiento del estado del agente.
    \item Integración de persistencia mediante Git en el sandbox, permitiendo al agente mantener un historial de cambios y facilitar la recuperación de estados previos.
    \item Optimización del sistema de prompts para modo agéntico, mejorando la capacidad de razonamiento y toma de decisiones del agente.
    \item Mejora de compatibilidad entre diferentes modelos LLM que ya implementan razonamiento previo, ampliando la flexibilidad del sistema.
\end{itemize}
\subsection{Sprint 5: Investigación e Integración de Hardware de Alto Rendimiento}

\textbf{Estado:} \textcolor{late}{\textbf{Completado (con extensión planificada)}}

\textbf{Fecha de inicio:} 09 de enero de 2026.
\textbf{Fecha de cierre:} 18 de enero de 2026.

\textbf{Comentarios:}
\begin{itemize}
    \item \textbf{Hito Crítico:} Implementación de sistema de gestión de ramas de conversaciones.
    \item Integración de persistencia mediante Git en el sandbox (solución técnica compleja y elegante).
    \item Acceso y pruebas con infraestructura Computaex de alto rendimiento.
    \item Optimización del sistema de prompts para modo agéntico.
    \item Mejora de compatibilidad entre diferentes modelos LLM que ya implementan razonamiento previo.
    \item \textit{Nota:} Esta extensión se justifica por la complejidad técnica y el valor añadido de la persistencia basada en Git.
\end{itemize}

\subsection{Sprint 6: Validación y Pruebas Intensivas}

\textbf{Estado:} \textcolor{early}{\textbf{Completado (acelerado)}}

\textbf{Fecha de inicio:} 19 de enero de 2026.
\textbf{Fecha de cierre:} 22 de enero de 2026.

\textbf{Comentarios:}
\begin{itemize}
    \item Ejecución intensiva de auditorías end-to-end automatizadas.
    \item Validación del sistema completo en entornos diversos.
    \item Pruebas de estrés y estabilidad con múltiples modelos.
    \item Análisis y corrección de falsos positivos.
    \item Sprint acelerado debido a la madurez del sistema tras las mejoras del Sprint 5.
\end{itemize}

\subsection{Sprint 7: Evaluación y Documentación Final}

\textbf{Estado:} \textcolor{inprogress}{\textbf{En progreso}}

\textbf{Fecha de inicio:} 23 de enero de 2026.
\textbf{Fecha estimada de cierre:} 04 de febrero de 2026.

\textbf{Comentarios:}
\begin{itemize}
    \item Redacción del borrador de la memoria (23-31 enero).
    \item Fase de revisión y refinamiento (01-04 febrero).
    \item Generación de gráficos, tablas de resultados y bibliografía.
    \item Preparación de la presentación y defensa.
\end{itemize}

% ============================================================
\subsection{Diagrama de Gantt: Ejecución Real Completa}

El siguiente diagrama muestra todo el proyecto completado con fechas reales del calendario.

\begin{figure}[H]
\centering
\begin{ganttchart}[
    hgrid,
    vgrid,
    x unit=0.18cm,
    y unit chart=0.6cm,
    bar height=0.5,
    title label font=\bfseries\tiny,
    bar label font=\tiny,
    time slot format=isodate,
    time slot unit=day
]{2025-11-17}{2026-02-04}

% Títulos manuales en español
\gantttitle{2025-2026}{80} \\
\gantttitle{Nov}{14}\gantttitle{Dic}{31}\gantttitle{Ene}{31}\gantttitle{Feb}{4} \\
\gantttitle{17-23}{7}\gantttitle{24-30}{7}\gantttitle{1-7}{7}\gantttitle{8-14}{7}\gantttitle{15-21}{7}\gantttitle{22-31}{10}\gantttitle{1-7}{7}\gantttitle{8-14}{7}\gantttitle{15-21}{7}\gantttitle{22-31}{10}\gantttitle{1-4}{4} \\

% Sprint 1: 17/11 - 23/11 (a tiempo = azul)
\ganttbar[bar/.append style={fill=ontime, draw=none}]{S1: Cimentación}{2025-11-17}{2025-11-23} \\

% Sprint 2: 24/11 - 02/12 (anticipado = verde)
\ganttbar[bar/.append style={fill=early, draw=none}]{S2: Sandbox KaliLinux}{2025-11-24}{2025-12-02} \\

% Sprint 3: 03/12 - 25/12 (a tiempo = azul)
\ganttbar[bar/.append style={fill=ontime, draw=none}]{S3: Integración LLM}{2025-12-03}{2025-12-25} \\

% Sprint 4: 26/12 - 08/01 (a tiempo = azul)
\ganttbar[bar/.append style={fill=ontime, draw=none}]{S4: Desarrollo MVP}{2025-12-26}{2026-01-08} \\

% Sprint 5: 09/01 - 18/01 (extendido = rojo claro)
\ganttbar[bar/.append style={fill=late, draw=none}]{\textcolor{late}{S5: Hardware Alto Rend.}}{2026-01-09}{2026-01-18} \\

% Sprint 6: 19/01 - 22/01 (acelerado = verde)
\ganttbar[bar/.append style={fill=early, draw=none}]{S6: Validación}{2026-01-19}{2026-01-22} \\

% Sprint 7: 23/01 - 04/02 (en progreso = amarillo)
\ganttbar[bar/.append style={fill=inprogress, draw=none}]{S7: Documentación}{2026-01-23}{2026-02-04} \\

% Relaciones
\ganttlink{elem0}{elem1}
\ganttlink{elem1}{elem2}
\ganttlink{elem2}{elem3}
\ganttlink{elem3}{elem4}
\ganttlink{elem4}{elem5}
\ganttlink{elem5}{elem6}

\end{ganttchart}
\caption{Cronograma de ejecución real completo (actualizado a 29/01/2026).}

% Leyenda de colores
\vspace{0.5em}
\begin{center}
\begin{tabular}{llll}
\colorbox{ontime}{\rule{0.8cm}{0.35cm}} & A tiempo &
\colorbox{early}{\rule{0.8cm}{0.35cm}} & Anticipado \\[0.3em]
\colorbox{late}{\rule{0.8cm}{0.35cm}} & Retrasado &
\colorbox{inprogress}{\rule{0.8cm}{0.35cm}} & En progreso \\
\end{tabular}
\end{center}

\end{figure}

\bigskip

% ============================================================
\section{Gestión de Desviaciones y Ajustes de Cronograma}

Durante la ejecución del proyecto se produjeron varias desviaciones respecto al cronograma inicial, las cuales se gestionaron de manera proactiva para mantener la calidad del entregable final.

\subsection{Desviaciones Identificadas}

\begin{description}[style=nextline, leftmargin=1.5cm, font=\bfseries]

    \item[Adelanto en Sprint 2 (Positivo)]
    El desarrollo de la infraestructura de red aislada se completó 5 días antes de lo previsto, generando un margen de maniobra que se aprovechó para el Sprint 3.
    
    \item[Identificación de Oportunidad Crítica (Sprint 4-5)]
    Durante el desarrollo del MVP se identificó la oportunidad de integrar:
    \begin{itemize}[leftmargin=0.5cm, label=$\circ$]
        \item Sistema de gestión de ramas de conversaciones
        \item Persistencia mediante Git en el sandbox
        \item Acceso a infraestructura Computaex de alto rendimiento
    \end{itemize}
    
    \item[Extensión Técnica del Sprint 5 (Justificada)]
    Se extendió el cronograma una semana (del 12 al 18 de enero) debido a:
    \begin{itemize}[leftmargin=0.5cm, label=$\circ$]
        \item Complejidad técnica de la persistencia basada en Git
        \item Integración y pruebas con modelos locales en Computaex
        \item Optimización avanzada del sistema de prompts
    \end{itemize}

\end{description}

\subsection{Estrategias de Mitigación Aplicadas}

\begin{enumerate}
    \item \textbf{Aprovechamiento de márgenes:} El adelanto del Sprint 2 proporcionó flexibilidad para absorber la complejidad técnica posterior.
    
    \item \textbf{Replanificación proactiva:} Se reestructuró el Sprint 5 como "Investigación e Integración de Hardware" en lugar de simplemente "Validación".
    
    \item \textbf{Aceleración estratégica:} El Sprint 6 se ejecutó de forma intensiva (4 días) aprovechando la madurez del sistema.
    
    \item \textbf{Justificación técnica:} Las desviaciones se fundamentaron en mejoras técnicas significativas que aportan valor al proyecto final.
\end{enumerate}

\subsection{Impacto en el Resultado Final}

La gestión de estas desviaciones resultó en:
\begin{itemize}
    \item \textbf{Mejora técnica sustancial:} Sistema más robusto con persistencia Git y compatibilidad multi-modelo.
    \item \textbf{Validación en hardware de alto rendimiento:} Pruebas con infraestructura Computaex.
    \item \textbf{Retraso controlado:} Únicamente una semana de desviación respecto al cronograma original.
    \item \textbf{Calidad del entregable:} Sistema final más completo y técnicamente avanzado.
\end{itemize}

\textbf{Conclusión:} Las desviaciones, aunque generaron un retraso menor, se tradujeron en mejoras técnicas significativas que justifican plenamente el ajuste temporal realizado.

Estas entradas documentan el avance real y facilitan la comparación posterior entre la planificación inicial y el desarrollo efectivo del proyecto.

\end{document}